\documentclass[aspectratio=169]{beamer}
\usetheme{Singapore}
\usecolortheme[RGB={0,110,110}]{structure}

\title[PRS in Cardiology]{Journal Club: Polygenic Risk Scores in Preventive Cardiology}
\subtitle{Appraisal of the HEARTLINE-2 cohort study}
\author[H. Lieu]{Hannah Lieu, MD}
\institute[St.\ Aldric]{Cardiology Fellowship Program, St.\ Aldric Medical Center}
\date{22 July 2026}

\begin{document}

\begin{frame}
  \titlepage
\end{frame}

\begin{frame}{Outline}
  \tableofcontents
\end{frame}

\section{Study design}

\begin{frame}{Study at a glance}
  \begin{itemize}
    \item Prospective cohort of 41,206 adults aged 40 to 69, median follow-up 8.2 years.
    \item Exposure: polygenic risk score for coronary artery disease, top quintile versus middle.
    \item Primary outcome: first major adverse cardiovascular event.
    \item Adjusted for pooled cohort equation risk, statin use, and family history.
  \end{itemize}
\end{frame}

\section{Results and critique}

\begin{frame}{Results and our critique}
  Top-quintile score carried a hazard ratio of 1.62 (95\% CI 1.44 to 1.83)
  after adjustment for clinical risk.
  \begin{itemize}
    \item Strength: event ascertainment via linked registries, minimal loss to follow-up.
    \item Concern: score derived from a European-ancestry panel, limiting transportability.
    \item Concern: no reclassification analysis, so clinical utility remains unclear.
  \end{itemize}
\end{frame}

\begin{frame}{Take-home}
  \begin{center}
    {\Large Promising for risk enrichment, not ready to change statin thresholds.}\\[1em]
    Next journal club: lipoprotein(a) screening trials, same time in two weeks.
  \end{center}
\end{frame}

\end{document}
